% =======================================================================================
\section{The Galactic-Binary Signal and the LISA Response}
\label{app:signal}
% =======================================================================================

We compute the signal with JaxGB \citep{jaxgb}. This section reproduces its model in its own
notation and conventions, which are also those of the time-domain response simulator LISA GW
Response \citep{lisagwresponse}: the two codes share the triad, the polarisation convention
and the link response, and pass the same eight parameters to each other unchanged. Only
the heterodyned evaluation (\cref{app:signal-heterodyne}) and the TDI combination
(\cref{app:signal-tdi}) are specific to JaxGB.

\subsection{Waveform}
\label{app:signal-waveform}

A detached compact binary in the Galaxy, most often a pair of white dwarfs, emits a
quasi-monochromatic gravitational wave at twice its orbital frequency
\citep{nelemans2001gravitational,korol2017prospects}, which drifts slowly through radiation
reaction. The transverse-traceless strain of the plane wave is
\begin{equation}
  h_{ij}(\xi) = h_+(\xi)\,e^+_{ij} + h_\times(\xi)\,e^\times_{ij},\qquad
  \xi = t - \frac{\hat{\mathbf k}\cdot\mathbf x(t)}{c},
  \label{eq:xi}
\end{equation}
with $\hat{\mathbf k}$ the direction of propagation, $\mathbf x$ the position and
$e^{+,\times}_{ij}$ the polarisation tensors of the source frame. For a mildly chirping,
circular binary the two polarisations are
\begin{equation}
  h_+(\xi) = -\Amp\,\bigl(1+\cos^2\iota\bigr)\cos\Psi(\xi),\qquad
  h_\times(\xi) = -2\Amp\cos\iota\,\sin\Psi(\xi),\qquad
  \Psi(\xi) = -\phi_0 + 2\pi f_0\,\xi + \pi\fdot\,\xi^2,
  \label{eq:waveform}
\end{equation}
with strain amplitude $\Amp$, inclination $\iota$ of the orbital angular momentum to the line
of sight, initial phase $\phi_0$, frequency $f_0$ and frequency derivative $\fdot$, all
referred to the start of the observation. The sky position, right ascension $\alpha$ and
declination $\delta$ in the equatorial (ICRS) frame of the orbit files, defines the
orthonormal triad
\begin{equation}
  \hat{\mathbf u} = (\sin\alpha,\,-\cos\alpha,\,0),\qquad
  \hat{\mathbf v} = (-\sin\delta\cos\alpha,\,-\sin\delta\sin\alpha,\,\cos\delta),\qquad
  \hat{\mathbf k} = (-\cos\delta\cos\alpha,\,-\cos\delta\sin\alpha,\,-\sin\delta),
  \label{eq:triad}
\end{equation}
and with it the polarisation tensors of the solar-system barycentre (SSB),
$\epsilon^+_{ij} = (\hat{\mathbf u}\otimes\hat{\mathbf u} - \hat{\mathbf v}\otimes\hat{\mathbf v})_{ij}$
and
$\epsilon^\times_{ij} = (\hat{\mathbf u}\otimes\hat{\mathbf v} + \hat{\mathbf v}\otimes\hat{\mathbf u})_{ij}$.
The source frame is rotated from the SSB frame by the polarisation angle $\psi$,
$e^+_{ij} = \cos2\psi\,\epsilon^+_{ij} + \sin2\psi\,\epsilon^\times_{ij}$ and
$e^\times_{ij} = -\sin2\psi\,\epsilon^+_{ij} + \cos2\psi\,\epsilon^\times_{ij}$, so that
\begin{equation}
  h_{ij}(\xi) = h_+^{\rm SSB}(\xi)\,\epsilon^+_{ij} + h_\times^{\rm SSB}(\xi)\,\epsilon^\times_{ij},
  \qquad
  h_+^{\rm SSB} = h_+\cos2\psi - h_\times\sin2\psi,\qquad
  h_\times^{\rm SSB} = h_+\sin2\psi + h_\times\cos2\psi .
  \label{eq:pol-ssb}
\end{equation}
In complex form, $h_{ij}(\xi) = \mathrm{Re}\bigl\{\bigl(A_+^{\rm SSB}\epsilon^+_{ij}
+ A_\times^{\rm SSB}\epsilon^\times_{ij}\bigr)e^{\ii\Psi(\xi)}\bigr\}$ with
\begin{equation}
  A_+^{\rm SSB} = -\Amp\bigl[(1+\cos^2\iota)\cos2\psi + 2\ii\cos\iota\sin2\psi\bigr],\qquad
  A_\times^{\rm SSB} = -\Amp\bigl[(1+\cos^2\iota)\sin2\psi - 2\ii\cos\iota\cos2\psi\bigr].
  \label{eq:complex-amp}
\end{equation}
In LISA GW Response the same triad reads $\hat{\mathbf k} = -\mathbf e_r$,
$\hat{\mathbf u} = -\mathbf e_\phi$, $\hat{\mathbf v} = -\mathbf e_\theta$ in the spherical basis
of the ICRS, and its galactic-binary strain is \cref{eq:waveform,eq:pol-ssb} term by term.

For a circular binary driven by gravitational radiation alone, the frequency derivative is
fixed by the chirp mass $\Mc$,
\begin{equation}
  \fdot = \frac{96}{5}\,\pi^{8/3}\left(\frac{G\Mc}{c^3}\right)^{5/3} f_0^{11/3},
  \label{eq:fdot}
\end{equation}
and the amplitude by the distance $D$, $\Amp = 2(G\Mc)^{5/3}(\pi f_0)^{2/3}/(c^4 D)$. We
parametrise each source by $\theta = (f_0, \Mc, \Amp, \alpha, \delta, \psi, \iota, \phi_0)$,
derive $\fdot$ from \cref{eq:fdot}, and do not sample the distance: the amplitude prior is
set by the signal-to-noise ratio instead (\cref{app:noise-snr}). Tides and mass transfer,
which change $\fdot$ for some systems, are neglected.

\subsection{Single-Link Response}
\label{app:signal-response}

LISA measures the wave as relative frequency fluctuations of six laser links between three
spacecraft. Let $i$ be the receiving and $j$ the emitting spacecraft of link $ij$, $t_i$ the
time of reception and $t_j = t_i - L_{ij}(t_i)$ the time of emission, with $L_{ij}$ the light
travel time. The link response is \citep{finn2009response,cornish2009alternative}
\begin{equation}
  y_{ij}(t_i) = -\frac12\,\frac{H_{ij}(\xi_i) - H_{ij}(\xi_j)}{1 - \hat{\mathbf k}\cdot\hat{\mathbf x}_{ij}},
  \qquad
  H_{ij}(\xi) = \hat x^k_{ij}\hat x^l_{ij}\,h_{kl}(\xi)
  = F^+_{ij}\,h_+^{\rm SSB}(\xi) + F^\times_{ij}\,h_\times^{\rm SSB}(\xi),
  \label{eq:link-response}
\end{equation}
with $\xi_i = t_i - \hat{\mathbf k}\cdot\mathbf x_i(t_i)/c$,
$\xi_j = t_j - \hat{\mathbf k}\cdot\mathbf x_j(t_j)/c$, the unit vector
$\hat{\mathbf x}_{ij} = (\mathbf x_i(t_i) - \mathbf x_j(t_j))/\lVert\mathbf x_i(t_i) - \mathbf x_j(t_j)\rVert$
from the emitter to the receiver, and the antenna patterns
$F^+_{ij} = (\hat{\mathbf u}\cdot\hat{\mathbf x}_{ij})^2 - (\hat{\mathbf v}\cdot\hat{\mathbf x}_{ij})^2$
and $F^\times_{ij} = 2(\hat{\mathbf u}\cdot\hat{\mathbf x}_{ij})(\hat{\mathbf v}\cdot\hat{\mathbf x}_{ij})$.
This is the response of LISA GW Response, $y_{ij}(t) = [H_{ij}(t - L_{ij} - \hat{\mathbf k}\cdot\mathbf x_j/c)
- H_{ij}(t - \hat{\mathbf k}\cdot\mathbf x_i/c)]/[2(1 - \hat{\mathbf k}\cdot\mathbf n_{ij})]$,
with $\mathbf n_{ij}\equiv\hat{\mathbf x}_{ij}$.

For the galactic binary, define
$\mathcal A_{ij} = A_+^{\rm SSB}F^+_{ij} + A_\times^{\rm SSB}F^\times_{ij}$ and
\begin{equation}
  \alpha_{ij}(t_i) = \pi L_{ij}(t_i)\,\bigl(1 - \hat{\mathbf k}\cdot\hat{\mathbf x}_{ij}\bigr)\,
  \bigl(f_0 + \fdot\,\xi_i\bigr).
  \label{eq:alpha}
\end{equation}
With $L_{ij}\simeq\lVert\mathbf x_i(t_i) - \mathbf x_j(t_j)\rVert/c$ and the adiabatic
approximation $\fdot(\xi_i - \xi_j)/2\ll f_0$ (no significant change of frequency during a
light travel time), half the phase difference between emission and reception is
$\alpha_{ij}$, and the response becomes
\begin{equation}
  y_{ij}(t_i) = -\mathrm{Re}\left\{\frac{\ii\,\mathcal A_{ij}\sin\alpha_{ij}}
  {1 - \hat{\mathbf k}\cdot\hat{\mathbf x}_{ij}}\,
  e^{-\ii\left(\alpha_{ij} - \Psi(\xi_i)\right)}\right\}.
  \label{eq:link-gb}
\end{equation}

\subsection{Heterodyned Frequency-Domain Response}
\label{app:signal-heterodyne}

With $q = \operatorname{round}(f_0\Tobs)$ the frequency bin closest to $f_0$, the response
separates into a fast carrier and a slowly varying envelope,
$y_{ij}(t_i) = \mathrm{Re}\bigl\{\mathcal G_{ij}(t_i)\,e^{2\pi\ii qt_i/\Tobs}\bigr\}$, with
\begin{equation}
  \mathcal G_{ij}(t_i) = -\frac{\ii}{2}\,\frac{f_0 + \fdot\,\xi_i}{f_{\star}}\,\mathcal A_{ij}\,
  \operatorname{sinc}\alpha_{ij}\;
  \exp\Bigl\{-\ii\Bigl[\alpha_{ij} + 2\pi f_0\frac{\hat{\mathbf k}\cdot\mathbf x_i(t_i)}{c}
  - \Bigl(-\phi_0 + 2\pi\bigl(f_0 - \tfrac{q}{\Tobs}\bigr)t_i + \pi\fdot\,\xi_i^2\Bigr)\Bigr]\Bigr\},
  \label{eq:slow}
\end{equation}
where $\operatorname{sinc}x = \sin x/x$ and $f_\star = (2\pi L)^{-1}$, with $L$ in seconds.
The envelope varies on the scale of LISA's orbit, so its spectrum is confined to a few bins
around zero: the Doppler sideband $\pm f_0v_\oplus/c$ ($v_\oplus = 29.8$\,km/s, so
$\sim10^{-4}f_0$) plus the chirp drift $\fdot\,\Tobs$. JaxGB samples it at $R$ times
$t_m = m\Tobs/R$ and Fourier-transforms it, which gives the one-sided spectrum on the $R$ bins
starting at $k_{\min} = q - R/2$:
\begin{equation}
  \tilde y_{ij}\bigl(f_{q+k}\bigr) = \frac{\Tobs}{2R}\sum_{m=0}^{R-1}\mathcal G_{ij}(t_m)\,
  e^{-2\pi\ii km/R},\qquad k = -\tfrac R2,\dots,\tfrac R2 - 1,\qquad f_k = \frac{k}{\Tobs}.
  \label{eq:heterodyne}
\end{equation}
The factor $1/2$ is the positive-frequency half of the real signal, and the normalisation is
that of the continuous transform, $\tilde h(f) \simeq \int_0^{\Tobs}h(t)\,e^{-2\pi\ii ft}\,\dd t$.
This is the fast galactic-binary response of \citet{cornish2007tests}, vectorised over
sources and differentiable. The baseline JaxGB model makes two further approximations,
accurate at the percent level: the emitter's position is taken at the time of reception,
$\mathbf x_j(t_j)\simeq\mathbf x_j(t_i)$, so that $\hat{\mathbf x}_{ji} = -\hat{\mathbf x}_{ij}$
and $\mathcal A_{ji} = \mathcal A_{ij}$, and all six light travel times are equal and constant,
$L_{ij} = L$, their average. With the equal-arm orbits of LISA Orbits \citep{lisaorbits},
$L = 2.5\,\text{Gm}/c = 8.34$\,s and $f_\star = 19.1$\,mHz.

\subsection{Time-Delay Interferometry}
\label{app:signal-tdi}

Laser frequency noise is cancelled by time-delay interferometry
\citep[TDI;][]{tinto2021time}. JaxGB forms the first-generation (1.5) Michelson
combinations for equal and constant arms on the envelopes of \cref{eq:slow}, with the
instantaneous frequency $f = f_0 + \fdot\,t$ and the delay phase $x = 2\pi fL$:
\begin{equation}
  X = -2\ii\sin x\;e^{-\ii x}\Bigl[\bigl(\mathcal G_{21} - \mathcal G_{31}\bigr)
  + e^{-\ii x}\bigl(\mathcal G_{12} - \mathcal G_{13}\bigr)\Bigr],
  \label{eq:tdi-x}
\end{equation}
with $Y$ and $Z$ from the cyclic permutation $1\to2\to3\to1$, Fourier-transformed as in
\cref{eq:heterodyne}. Since $-2\ii\sin x\,e^{-\ii x} = -(1 - D^2)$ with $D = e^{-\ii x}$ the
delay operator, \cref{eq:tdi-x} is $(1-D^2)$ times a combination of the four links that
involve spacecraft 1. The standard equal-arm TDI-1.5 Michelson combination
\citep[eq.~14]{babak2021lisa} pairs the delays the other way round,
$(1-D^2)\bigl[(y_{13}-y_{12}) + D\,(y_{31}-y_{21})\bigr]$; the two differ by
$(1-D^2)(1-D)\bigl[(y_{13}-y_{31}) - (y_{12}-y_{21})\bigr]$, which is of higher order in $x$.
On JaxGB's own link responses, for sources drawn from our prior, they agree to
$2.4\times10^{-4}$ (relative) up to 3\,mHz and to $7\times10^{-3}$ at 10\,mHz, with powers
equal to 0.2\%. We use the quasi-orthogonal channels \citep{prince2002lisa}
\begin{equation}
  A = \frac{Z-X}{\sqrt2},\qquad E = \frac{X-2Y+Z}{\sqrt6},\qquad T = \frac{X+Y+Z}{\sqrt3},
  \label{eq:aet}
\end{equation}
whose noises are uncorrelated for equal arms (\cref{app:noise-psd}). The isotropic sky prior
does not depend on the frame; sources quoted in ecliptic coordinates, as in the LISA Data
Challenges, must be rotated to the equatorial frame of the orbits.

\subsection{Frequency Windows}
\label{app:signal-windows}

A training example is one window of $K=4R$ consecutive frequency bins, starting at an even
bin $k_0$, with centre frequency $\fw = (k_0 + K/2)/\Tobs$. The window holds $S=4$ sources,
whose $R$-bin responses are added at their own offsets $k_{\min}-k_0$. A guard band of $R/2$
bins at each edge keeps every response inside the window: sources are drawn with
\begin{equation}
  f_0 \in \bigl[f_{\rm lo}(w), f_{\rm hi}(w)\bigr]
  = \Bigl[\max\bigl(\tfrac{k_0 + R/2}{\Tobs}, f_{\min}\bigr),\;
          \min\bigl(\tfrac{k_0 + K - R/2}{\Tobs}, f_{\max}\bigr)\Bigr],
  \label{eq:window-interior}
\end{equation}
an interior of $3R$ bins. The window is placed by drawing $f_c$ log-uniformly on
$[f_{\min}, f_{\max}]$ and snapping the start to the nearest even bin
$k_0 = 2\operatorname{round}(f_c\Tobs/2 - R)$, clipped to the valid range. Windows therefore
tile the band with density $\propto 1/f$, and $f_0$ is log-uniform over the band to within
the snapping.

\paragraph{How many bins a source needs.}
Half of a source's response window must hold the Doppler sideband plus the whole chirp
drift, so a source needs
\begin{equation}
  n_{\rm span}(f_0,\Mc) = 2\left\lceil\Bigl(\frac{f_0 v_\oplus}{c}
  + \fdot(\Mc, f_0)\,\Tobs\Bigr)\Tobs\right\rceil \le R
  \label{eq:span}
\end{equation}
bins. For a given $R$, the highest frequency the prior can reach is set by the heaviest
binary, $\Mc_{\max} = 1.4\cdot2^{-1/5} = 1.219\,\Msun$, and solves
$(v_\oplus/c)\,f_0 + \kappa\,\Mc_{\max}^{5/3}\Tobs f_0^{11/3} = R/(2\Tobs)$, with
$\kappa$ the constant of \cref{eq:fdot}. In $u = f_0^{1/3}$ this
is a degree-11 trinomial, solved numerically by Newton's method from the smaller of the
Doppler-only root $Rc/(2\Tobs v_\oplus)$ and the chirp-only root. The two terms scale very
differently: the Doppler term is linear in $f_0$ and the chirp term goes as $f_0^{11/3}$, so
once the chirp dominates (above $\sim4$\,mHz for the heaviest binaries,
\cref{fig:band}c) quadrupling $R$ buys only $4^{3/11}\simeq1.46$ times more bandwidth.

\paragraph{Rungs.}
The problem comes in three sizes that differ only in $R$ and hence in how much of the band a
window covers (\cref{tab:rungs}). With the full band, $R$ is sized automatically as the
power of two above $n_{\rm span}(f_{\max},\Mc_{\max})$, with a floor of 256 points: at
12\,mHz the Doppler sideband is 75 bins and the chirp of the heaviest binary 291 bins, so
$n_{\rm span}=734$ and $R=1024$, which could hold the heaviest binaries up to 13.3\,mHz.

\begin{table}[h]
\caption{The three problem sizes. They share the priors, the four sources per window, the
network and the training recipe; only the response points $R$, and with them the window and
the number of data tokens, differ. The band fraction is of $\log f$ over 0.1--12\,mHz.}
\label{tab:rungs}
\centering
\small
\begin{tabular}{lrrrrll}
\toprule
rung & $R$ & window $K$ [bins] & interior [bins] & data tokens & $f_0$ range & band \\
\midrule
XS & 16 & 64 & 48 & 32 & 0.1--1.26\,mHz & 53\% \\
S & 64 & 256 & 192 & 128 & 0.1--4.15\,mHz & 78\% \\
B & 1024 & 4096 & 3072 & 2048 & 0.1--12\,mHz & 100\% \\
\bottomrule
\end{tabular}
\end{table}

\begin{figure}[h]
  \centering
  \includegraphics[width=\textwidth]{figures/band.pdf}
  \caption{The band as the training prior sees it. (a) The TDI-1.5 noise PSDs of the $A$,
  $E$ and $T$ channels in fractional-frequency units, which whiten the data
  (\cref{eq:psd}). (b) The amplitude prior: the band between a sky-averaged SNR of 7 and of
  1000 at each frequency (\cref{eq:amp-window}); the dot is the faint source of
  \citet{strub2022bayesian} at SNR 7.93, used to validate the likelihood (\cref{app:mcmc}).
  (c) The half-span of a source's response, Doppler sideband plus chirp drift over
  $\Tobs$ (\cref{eq:span}), against the half-window of each rung; the dots mark the highest
  frequency each rung can hold for the heaviest binary in the prior: 1.26, 4.15 and
  13.3\,mHz (the prior itself stops at 12\,mHz).}
  \label{fig:band}
\end{figure}

\subsection{What the Simulator Leaves Out}
\label{app:signal-limits}

The training distribution is deliberately simpler than real LISA data. Every window holds
exactly four resolvable binaries and nothing else: there is no confusion foreground from the
millions of unresolved binaries \citep{timpano2006characterizing,karnesis2021characterization},
no other source class, no gaps or glitches, and the noise is stationary and known. Sources in
neighbouring windows never leak in, because the guard band keeps every response inside its
own window. For $R<256$ (XS and S) the slow response is sampled critically rather than
generously, which leaves per-bin amplitudes of the highest-frequency XS sources with a
$\sim$10--25\% aliasing error against a 1024-point reference while preserving the total
power; the same simulator generates the training data and the test injections, so this
affects realism, not consistency. Using the posteriors inside a global fit would require
lifting the first two simplifications: a variable number of sources per window, and sources
that straddle window edges.
